\documentclass[10pt,twocolumn,letterpaper]{article}

\usepackage{times}
\usepackage{epsfig}
\usepackage{graphicx}
\usepackage{amsmath}
\usepackage{amssymb}
\usepackage{booktabs}
\usepackage{hyperref}

\title{VLM-Diff: Deterministic-First Description of UI Visual Regressions\\with VLM Escalation Only Where the DOM Cannot Explain}

\author{
Anonymous Author(s) \\
Institution \\
{\tt\small email@domain.com}
}

\begin{document}
\maketitle

\begin{abstract}
Frontier vision-language models (VLMs) struggle with fine-grained visual differences, achieving only 62--78\% accuracy on subtle UI changes vs a 95\% human baseline. We propose a hybrid approach that leverages the UI domain's unique advantage: DOM structure as ground truth. Our pipeline localizes changed regions with deterministic DOM+pixel diffing (no-change suppression: zero DOM changes $\Rightarrow$ unchanged, 0\% false positives on all 6 no-change fixtures of a 145-pair dataset), then describes each region on a \textit{tiered} basis: regions the DOM diff fully explains (color, text, style, lifecycle, geometry) receive template-generated descriptions with zero tokens; only regions with no DOM signal (e.g.\ canvas repaints) escalate to a VLM. Real-API MVPs (Kimi K3) confirm the architecture: a three-arm run shows +33pp recall over full-image VLM with 0\% false positives and refutes plain crop-then-classify (58--67\% vs 100\% conditional accuracy) — passing the detector's changed fields as a text hint recovers the gap (83.3\%). A four-arm run then validates the tiered pipeline end-to-end: \textbf{100\% change-type accuracy at ${\sim}$15 tokens/pair} (74/75 regions deterministic, 98.9\% token savings), fixing both failure modes the VLM arm made — and a cross-vendor replication on qwen3.8-max reproduces the result exactly (100\%/0\%/100\%, same 1.3\% escalation). Offline replay of the tiered router over all 145 pairs shows \textbf{99.3\% of regions are fully deterministic}, with root-cause-first pair-level typing at \textbf{90.6\% end-to-end change-type accuracy at zero VLM tokens} (vs 82.7\% for largest-region-first typing). Code and dataset: \url{https://github.com/shidesheng0218/vlm-diff}
\end{abstract}

\section{Introduction}

Visual regression testing—detecting unintended UI changes in web applications—is critical for continuous deployment pipelines. Traditional tools (Percy, Applitools, Chromatic) rely on pixel-diff algorithms that produce high false-positive rates from anti-aliasing noise, font rendering jitter, and dynamic content~\cite{applitools2026}.

Recent work has explored using vision-language models (VLMs) for fine-grained image comparison. However, VLM-SubtleBench~\cite{vlmsubtlebench2026} showed that even frontier models struggle: GPT-5-thinking achieved only 77.8\% accuracy on near-identical image pairs (vs 95.5\% human baseline), with spatial shifts performing worst at 59.9\%. Concatenating before/after into one image hurt 9 of 10 categories. OmniDiff~\cite{omnidiff2026} demonstrated a 6$\times$ performance gap between fine-tuned specialists (CIDEr 31.3) and zero-shot GPT-4o (5.2).

\textbf{Key insight}: UI screenshots have a unique advantage over generic images—\textit{DOM structure as ground truth}. This advantage is stronger than mere change detection: the DOM diff does not just say \textit{that} something changed, it records \textit{which computed properties} changed and their before/after values. A color delta from \texttt{rgb(37,99,235)} to \texttt{rgb(220,38,38)} is fully explained by the DOM; asking a VLM to re-derive it from crops is redundant work that costs tokens and adds variance.

We therefore propose VLM-Diff, a tiered hybrid pipeline:
\begin{enumerate}
\item \textbf{Deterministic detection}: DOM diff + perceptual pixel diff fusion localizes candidate regions. If DOM reports zero changes, the pair is classified as unchanged \textit{regardless of pixel noise}.
\item \textbf{Deterministic-first description}: regions whose DOM evidence fully explains the change receive template-generated descriptions with zero tokens.
\item \textbf{VLM escalation}: only regions with no DOM signal (pixel-only repaints) are cropped and classified by a VLM.
\end{enumerate}

\textbf{Contributions}:
\begin{itemize}
\item A UI regression benchmark with DOM snapshots + ground-truth rects (145 pairs, 6 fixtures, 24 mutation types), and a deterministic detection layer with 100\% recall on changed pairs and 0\% FP rate (confirmed)
\item A real-API three-arm MVP validating the architecture (+33pp recall, 0\% FP, 0\% run-to-run variance on detection) and \textit{refuting} naive crop-then-classify, while showing a DOM-field text hint recovers the gap
\item A deterministic-first tiered description architecture in which 99.3\% of detected regions require no VLM call, achieving 90.6\% end-to-end change-type accuracy at zero token cost — with a root-cause attribution rule for pair-level typing in layout cascades
\item Open-source implementation, evaluation framework, and a zero-cost CI regression gate on the escalation rate
\end{itemize}

\section{Related Work}

\textbf{VLM benchmarks for subtle differences}. VLM-SubtleBench~\cite{vlmsubtlebench2026} evaluated 13K near-identical image pairs across 10 categories. Best model (GPT-5-thinking) achieved 77.8\%, with spatial/temporal categories $<$70\%. OmniDiff~\cite{omnidiff2026} showed fine-tuned specialists outperform zero-shot frontier VLMs by 6$\times$. WUICC-bench~\cite{wuicc2026} released the first UI regression dataset but without DOM structure or frontier VLM evaluation.

\textbf{Visual regression tools}. Commercial tools (Percy, Chromatic, Applitools) use deterministic pixel diff + ML-based content pattern matching. Applitools publicly opposes ``LLM as comparator''~\cite{applitools2026mcp}, citing flakiness concerns. Our deterministic tier directly addresses the flakiness critique: 0\% of no-change pairs are flagged, byte-identically across runs. DiffShot-AI~\cite{diffshot2026} uses VLMs to analyze code changes and generate test screenshots—orthogonal to our work.

\textbf{VLM attention mechanisms}. Chain-of-Focus~\cite{chainoffocus2026} and SPARC~\cite{sparc2026} propose learned iterative zoom/attention for fine-grained tasks. Our approach uses \textit{rule-based} DOM diff for localization (zero-shot, UI-domain-specific) rather than learned attention — and goes further by making the structural signal sufficient to \textit{describe} most changes without any model call.

\section{Method}

\subsection{Dataset Generation}

We render 6 realistic UI fixtures (card grid, form, navbar, data table, modal dialog, dashboard) in headless Chromium and apply 24 programmatic mutations across 7 categories:

\begin{itemize}
\item \textbf{Spatial shift} (3/6/8/12/28px, horizontal + vertical, margin-first with transform fallback)
\item \textbf{Color change} (background, text, border; subtle hue shifts and high-contrast swaps)
\item \textbf{Size change} ($\times$0.9--1.35 via CSS \texttt{transform})
\item \textbf{Text change} (similar-length, different, shortened, numeric)
\item \textbf{Element add/remove} (first/last child)
\item \textbf{Style change} (font-weight, border-radius, box-shadow, opacity)
\item \textbf{No-change} (identical DOM, re-screenshot for AA noise)
\end{itemize}

Each pair includes \texttt{before.png}/\texttt{after.png} (960$\times$500), \texttt{domBefore}/\texttt{domAfter} (JSON: element path, tag, id, rect, text, 7 computed styles), \texttt{groundTruthRect} (bbox of the mutated element), and category ground truth. Total: \textbf{145 pairs} (139 changed, 6 no-change). Layout-cascade effects (siblings displaced by an add/remove/move) are preserved: the DOM snapshot records every displaced element's rect, and ground-truth rects point at the mutated element, not the cascade.

\subsection{Stage 1: Deterministic Detection}

\textbf{DOM diff}. Each element's structural path, bounding rect, text, and computed styles are compared node-by-node between snapshots. Sub-pixel jitter ($<$1px) is tolerated; rect deltas are decomposed into \texttt{position} (signed dx/dy) vs \texttt{size} (signed dw/dh) so translation is never confused with scaling. Every change records its before/after property values — the substrate for tiered description.

\textbf{Perceptual diff}. pixelmatch~\cite{pixelmatch} (threshold 0.15, \texttt{includeAA:false}) with connected-component grouping into candidate rects.

\textbf{Fusion}. \textit{Zero DOM changes $\Rightarrow$ \texttt{changed:false}} — the no-change suppression rule. On all 139 changed pairs the detector fires (100\% recall) and on all 6 no-change pairs it stays silent (0\% FP), byte-identically across repeated runs. DOM change rects become primary candidates (corroborated by overlapping pixel regions where present); pixel-only regions are appended as fallback candidates.

\subsection{Stage 2: Tiered Description (Deterministic-First)}

The core contribution. Each candidate region is routed and described from the DOM diff alone whenever the DOM fully explains the change:

\textbf{Templates}. Color deltas are named via a small perceptually-weighted anchor palette (``background color changed from blue (\texttt{rgb(37,99,235)}) to red (\texttt{rgb(220,38,38)})''). Text deltas phrase numeric changes as values and others as wording. Style properties (font-weight, border-radius, box-shadow, opacity, border) each have a template; element add/remove uses lifecycle sentences; geometry renders signed deltas (``moved 28px right'', ``resized by +35\%''). Deterministic descriptions carry confidence 1.0 by construction.

\textbf{Root-cause attribution}. Real layout changes cascade: removing a card shifts every sibling's rect. A geometry-only region (only \texttt{position}/\texttt{size} changed) in a pair that also contains a non-geometry change (text/color/style/lifecycle) is a \textit{reflow follower}, worded as ``moved 24px up as part of a layout shift caused by a nearby change''. Pair-level typing prefers root-cause regions over larger followers: in an element-remove pair whose container shrinks (a follower that would naively be typed \textit{size-change} and is typically the \textit{largest} region), the lifecycle region wins.

\textbf{Escalation policy}. The only escalation trigger is a \textit{pixel-only} region — a visual change with no DOM signal (canvas repaint, mid-frame animation). Escalated regions are cropped (16px padding) and classified by the VLM with the detector's changed fields as a text hint (\S\ref{sec:mvp}).

The router's conservatism is deliberate and honest about circularity: because every mutation in the dataset is DOM-observable by construction, near-total determinism is expected \textit{on this benchmark}. The contribution is the architecture plus a calibrated escalation-rate measurement (\S\ref{sec:replay}); a zero-cost CI gate (\texttt{replay:assert}) fails any change that pushes previously-deterministic regions to the VLM.

\subsection{Stage 3: VLM Classification (Escalation Path)}

For escalated regions, the VLM receives cropped before/after windows (never full concatenated images, per VLM-SubtleBench) and a system prompt with the DOM-field hint, and returns strict JSON: \texttt{\{changeType, description, confidence\}}. Classifications are content-hash cached (crop pixels + prompt) with a swappable store.

\section{Evaluation}

\subsection{Baselines and Arms}

\begin{itemize}
\item \textbf{rawPairToVlm}: full before/after images $\rightarrow$ VLM, ``what changed and where?''
\item \textbf{pixelDiffOnly}: perceptual diff only (no DOM, no VLM)
\item \textbf{fullPipeline}: fusion $\rightarrow$ per-region VLM classification, optionally with the DOM-field hint (ablation)
\item \textbf{tieredPipeline}: fusion $\rightarrow$ deterministic-first description $\rightarrow$ VLM only for pixel-only regions
\end{itemize}

\subsection{Detection Results (All 145 Pairs, No API)}

Detection is fully deterministic: \textbf{139/139 changed pairs detected, 0/6 false positives}. pixelDiffOnly flags 15--25\% of no-change pairs in published systems; DOM suppression eliminates this class entirely.

\subsection{Real-API MVP\label{sec:mvp}}

A stratified 15-pair subset (5 subtle/boundary, 7 clear, 3 no-change), three arms, Kimi K3 via DashScope, two independent runs (2026-08-19):

\begin{table}[h]
\centering
\small
\begin{tabular}{lccc}
\toprule
\textbf{Metric} & \textbf{raw} & \textbf{no hint} & \textbf{+ hint} \\
\midrule
Recall (12 changed) & 66.7 & \textbf{100} & \textbf{100} \\
FP rate (3 no-change) & 33.3 & \textbf{0} & \textbf{0} \\
Classification acc. & \textbf{100} & 58--67 & 83.3 \\
End-to-end type acc. & 67--75 & 67 & \textbf{83.3} \\
Avg input tok/pair & $\sim$1505 & $\sim$583 & $\sim$692 \\
\bottomrule
\end{tabular}
\caption{Two-run range (Kimi K3). rawPairToVlm missed 3--4 subtle pairs and produced 1 false positive (run 2); both pipeline arms were byte-identical across runs.}
\label{tab:mvp}
\end{table}

\textbf{Findings}. (1) \textit{Detection is reproducibly won by the pipeline}: +33pp recall, 0\% FP, and zero run-to-run variance, vs a raw arm that fluctuates between runs. (2) \textit{Plain crop-then-classify is refuted}: identical crops without the hint scored 58.3--66.7\% conditional accuracy vs raw's 100\% — cropping amputates the reference frame type judgments need. (3) \textit{The DOM-field hint recovers the gap} (+25pp/+16.7pp across runs, $\sim$50 tokens) and wins end-to-end outright; hinted outputs are also terser and steadier ($\sim$295 vs $\sim$560 output tokens/pair). Caveats: small sample (95\% CI $\approx\pm$20pp), one model, one vendor; replication on more pairs and Claude/GPT-class models is pending.

\subsection{Tiered Pipeline: Offline Replay\label{sec:replay}}

Replaying the router over all 145 pairs with zero API calls:

\begin{table}[h]
\centering
\small
\begin{tabular}{lrrrr}
\toprule
\textbf{Kind} & \textbf{regions} & \textbf{det.} & \textbf{vlm} & \textbf{esc.} \\
\midrule
size-change & 136 & 136 & 0 & 0\% \\
spatial-shift & 122 & 122 & 0 & 0\% \\
element-remove & 122 & 122 & 0 & 0\% \\
color-change & 94 & 94 & 0 & 0\% \\
text-change & 62 & 62 & 0 & 0\% \\
element-add & 41 & 41 & 0 & 0\% \\
style-change & 34 & 30 & 4 & 12\% \\
\midrule
\textbf{total} & \textbf{611} & \textbf{607} & \textbf{4} & \textbf{0.7\%} \\
\bottomrule
\end{tabular}
\caption{Region routing over 139 changed pairs. The only escalations are pixel-only repaint regions in three box-shadow pairs.}
\label{tab:replay}
\end{table}

Pair-level typing (offline, deterministic regions only): root-cause-first selection scores \textbf{90.6\%} (126/139) vs \textbf{82.7\%} for largest-region-first — and exceeds the full VLM pipeline's 83.3\% end-to-end accuracy on the MVP subset at \textit{zero} token cost. The gap between the two selection policies is precisely the cascade-attribution effect described in \S3.3: on element-add/remove and text-change pairs, the largest region is typically an innocent reflow follower.

Cost: the tiered pipeline's expected VLM spend on DOM-observable regressions is $\sim$0, vs $\sim$692 input tokens/pair for the hinted pipeline (and $\sim$1505 for raw) — a $>$99\% reduction on this benchmark, bounded by the (unknown, higher) escalation rate on real pages containing canvas and third-party widgets.

\subsection{Real-API Four-Arm Validation}

The tiered arm joined the live comparison: 36 pairs (30 changed, 6 no-change), Kimi K3 via DashScope, all four arms in one run (total cost \$0.18):

\begin{table}[h]
\centering
\small
\begin{tabular}{lcccc}
\toprule
\textbf{Metric} & \textbf{tiered} & \textbf{+ hint} & \textbf{no hint} & \textbf{raw} \\
\midrule
Recall (30 changed) & \textbf{100} & 100 & 100 & 76.7 \\
FP rate (6 no-change) & \textbf{0} & 0 & 0 & 0 \\
Type accuracy & \textbf{100} & 93.3 & 56.7 & 95.7$^{\dagger}$ \\
Tokens/pair (in+out) & \textbf{11+4} & 910+388 & 782+611 & 1505+196 \\
\bottomrule
\end{tabular}
\caption{Four-arm MVP (2026-08-28). $^{\dagger}$conditional on detection: raw missed 7 of 30 changed pairs (3 small color changes, 2 tiny spatial shifts) — the same failure class as both 15-pair runs.}
\label{tab:fourarm}
\end{table}

The tiered arm wins on every axis. Its two VLM-arm counterparts' errors are precisely the cases the deterministic tier eliminates by construction: a container's reflow growth out-voting an element-add at pair level (root-cause attribution fixes this), and a border-color change mis-typed as a style change (the template names the property). The no-hint refutation replicates a third time (56.7\%, after 58.3/66.7\%), and raw's recall ceiling is stable while its conditional accuracy stays high — the VLM fails at \textit{finding}, not at \textit{describing}. Caveats: one model, one vendor, one run; both escalation rates measured (0.7\% over 145 pairs offline, 1.3\% over this subset live) apply to DOM-observable mutations, so real-page escalation will be higher.

\subsection{Multi-Model Replication}

The four-arm protocol rerun on a second vendor's model (qwen3.8-max, DashScope):

\begin{table}[h]
\centering
\small
\begin{tabular}{lcccccc}
\toprule
 & \multicolumn{3}{c}{\textbf{tiered}} & \textbf{hint} & \textbf{no hint} & \textbf{raw} \\
\textbf{Model} & recall & FP & type-acc & type-acc & type-acc & recall \\
\midrule
Kimi K3 & 100 & 0 & \textbf{100} & 93.3 & 56.7 & 76.7 \\
qwen3.8-max & 100 & 0 & \textbf{100} & 90.0 & 90.0 & 83.3 \\
\bottomrule
\end{tabular}
\caption{Cross-vendor replication (36 pairs each). Escalation rate is identical (1.3\%, 1/75 regions) on both models — by construction, the deterministic tier's routing and descriptions are model-independent.}
\label{tab:multimodel}
\end{table}

Three findings replicate across vendors. (1) The tiered arm wins identically on both models — its routing and descriptions are byte-identical; only the single escalated region's classification depends on the model. (2) The raw arm's recall ceiling (76.7/83.3\%) is a cross-vendor failure mode: VLMs fail at \textit{finding} subtle changes, not at describing them. (3) New: the crop-then-classify weakness is \textit{model-dependent} — Kimi collapses to 56.7\% without the DOM hint while qwen3.8-max holds 90\%, so the reference-frame amputation effect varies in severity by model; on both models, however, the hint arm never beats the tiered arm, and the deterministic tier sidesteps the output-variance problem entirely (qwen's hinted arm averages 2321 output tokens/pair, 6$\times$ Kimi's).

\subsection{Blind Judge: Description Quality}

Type accuracy aside, is template text \textit{as good} as VLM text for a human triaging a regression? A blind judge (qwen3.8-max — a different vendor than the Kimi K3 that generated the VLM descriptions, avoiding self-preference bias) scored both arms on 30 pairs with labels randomized and ground truth provided:

\begin{table}[h]
\centering
\small
\begin{tabular}{lcc}
\toprule
\textbf{Dimension (1--5)} & \textbf{tiered} & \textbf{VLM} \\
\midrule
accuracy & 3.90 & \textbf{4.27} \\
specificity & \textbf{3.90} & 3.73 \\
readability & \textbf{4.60} & 4.50 \\
\midrule
\textbf{winner} & \textbf{14} & \textbf{14} (2 ties) \\
\bottomrule
\end{tabular}
\caption{Blind judge comparison. A split decision: templates win specificity by carrying exact values; the VLM wins accuracy via semantic naming (``the blue Confirm button was removed'' vs ``element \#btn-confirm was removed'').}
\label{tab:judge}
\end{table}

Judge noise is itself instructive: on one pair the judge rated a factually \textit{wrong} VLM description (a container narrowing; ground truth is element-add) 5/4/5 over the template's correct one — so exact-match type scoring (tiered 100\% vs 93.3\%) remains the more reliable yardstick, and judge scores should be read as comparative signal. The template improvement path is explicit: carry the element's visible text through the diff and lead with it. Comparable quality at zero marginal cost completes the tiered arm's value proposition.

\section{Discussion}

\textbf{Limitations}. (1) \textit{Escalation-rate circularity}: the benchmark's mutations are DOM-observable by construction, so near-total determinism is expected here; real-page escalation (canvas, images, third-party widgets) is the open empirical question. (2) Synthetic fixtures: 6 hand-built pages; no responsive breakpoints, virtualized tables, or live widgets. (3) The no-change FP denominator is 6 (95\% CI upper bound $\approx$39\%). (4) MVP validated on one mid-tier model and 15 pairs; judge-scored description quality for template vs VLM phrasing is pending (type accuracy needs no judge — it is exact-match against ground truth). (5) Chromium only.

\textbf{Future work}. (1) Blind LLM-judge scoring of deterministic vs VLM description phrasing, and four-arm replication on Claude/GPT-class models. (2) Real-page validation: two-commit checkout pairs of open-source repos with git-diff ground truth. (3) Severity classification (breaking vs cosmetic) for CI gating. (4) Fine-tuning a specialist on the deterministic descriptions themselves — the tier provides free training labels.

\section{Conclusion}

DOM structure is not just a detection signal for UI regression — it is a description signal. Fusing DOM+pixel diffs localizes changes perfectly on our benchmark (100\% recall, 0\% FP); routing each region through a deterministic-first describer with root-cause attribution answers 99.3\% of regions with zero tokens, escalating to a VLM only where no DOM signal exists. The live four-arm runs confirm it end-to-end on two vendors' models: 100\% change-type accuracy at ${\sim}$15 tokens/pair, with identical escalation rates; a blind cross-vendor judge rates the zero-token template descriptions on par with VLM descriptions (14/14/2). Where VLM calls remain necessary, the detector's changed fields are the difference between refuted (56--90\%, model-dependent) and winning (90--93\%) classification. Code, dataset, and a zero-cost CI gate on the escalation rate are open-sourced.

\bibliographystyle{plain}
\begin{thebibliography}{9}

\bibitem{vlmsubtlebench2026}
VLM-SubtleBench Authors.
\textit{VLM-SubtleBench: Evaluating Frontier Models on Near-Identical Image Pairs}.
arXiv:2603.07888, March 2026.

\bibitem{omnidiff2026}
OmniDiff Authors.
\textit{OmniDiff: Fine-Grained Image Difference Captioning at Scale}.
arXiv:2503.11093, March 2026.

\bibitem{wuicc2026}
WUICC-bench Authors.
\textit{WUICC-bench: A Benchmark for Web UI Visual Regression Testing}.
arXiv:2607.01728, July 2026.

\bibitem{applitools2026}
Applitools.
\textit{Applitools Eyes Documentation}.
\url{https://applitools.com/docs/eyes/playwright/core-concepts}, 2026.

\bibitem{applitools2026mcp}
Applitools.
\textit{Add Visual Testing to Your AI Workflow with the Applitools MCP Server}.
Blog post, January 2026.

\bibitem{diffshot2026}
DiffShot-AI.
\textit{AI-Powered Screenshot Testing}.
\url{https://github.com/sgasser/diffshot-ai}, 2026.

\bibitem{chainoffocus2026}
Chain-of-Focus Authors.
\textit{Chain-of-Focus: Adaptive Iterative Attention for VLMs}.
arXiv:2505.15436, May 2026.

\bibitem{sparc2026}
SPARC Authors.
\textit{SPARC: Visual Search and Reasoning for VLMs}.
arXiv:2602.06566, February 2026.

\bibitem{pixelmatch}
Mapbox.
\textit{pixelmatch: Smallest, simplest pixel-level image comparison library}.
\url{https://github.com/mapbox/pixelmatch}, 2024.

\end{thebibliography}

\end{document}
